\bookchapter{Глава 27: Родителям тоже будет легко — полное руководство по адаптации телефона для пожилых}

\chapterintro{Честно говоря, когда я писал эту главу, я думал о своих родителях. Каждый раз, когда на телефоне мамы всплывало уведомление «Обновление системы», она боялась нажать — вдруг сломает телефон. А шрифт в WeChat у папы был меньше муравья, но он не знал, как его увеличить — так и щурился два года. Это не значит, что они «глупые» — просто телефоны по умолчанию рассчитаны на людей около двадцати лет. В этой главе я самым простым языком разберу одну за другой самые болезненные проблемы, с которыми сталкиваются пожилые при использовании телефона.}

\booksection{Ноль. Сначала настроим правильный настрой: дело не в том, что «пожилые не разбираются», а в том, что «телефон сделан плохо»}

Прежде чем делать какие-либо настройки, нужно прояснить одну мысль:\textbf{интерфейсы современных смартфонов изначально почти не учитывали пожилых пользователей}. Всплывающие окна слишком маленькие и неразборчивые, жесты управления требуют точного касания, значки выглядят почти одинаково — это не проблема пользователя, это проблема продукта. Поэтому, помогая пожилым людям осваивать телефон, первое дело — не учить их, а помочь\textbf{настроить среду}. Как только среда настроена, большая часть проблем исчезнет.

Три железных правила при работе с пожилыми: (1)\textbf{Не торопите}— вы пользуетесь смартфоном десять лет, а они, возможно, всего три месяца. (2)\textbf{Используйте аналогии}— «Wi-Fi как водопровод дома: откроешь кран — потечёт вода». Такие сравнения понятны в десять тысяч раз лучше профессиональных терминов. (3)\textbf{Запишите}— устного объяснения может быть недостаточно, они могут забыть. Возьмите лист бумаги, напишите ключевые шаги и наклейте на обратную сторону чехла телефона — ничего действеннее.

\booksection{Один. Первый шаг: включите «Простой режим» — телефон станет чище и понятнее}

Практически все современные телефоны имеют встроенный «Простой режим» (иногда называется «Режим заботы», «Режим для пожилых», «Режим для старшего поколения»). После включения телефон кардинально преображается:

\begin{itemize}
  \item \textbf{Шрифт и значки становятся очень крупными}— больше не нужно надевать очки для чтения, чтобы разглядеть экран.
  \item \textbf{Значки на рабочем столе сильно упрощаются}— остаются только Телефон, Сообщения, Камера, Галерея, WeChat и несколько самых нужных приложений; всё остальное убирается.
  \item \textbf{Управление становится проще}— многие лишние функции (левая панель, умный ассистент, плавающие окна и т.д.) по умолчанию отключаются, интерфейс становится максимально чистым.
  \item \textbf{Звук становится громче}— мелодия входящего вызова, громкость уведомлений и звук в разговоре автоматически устанавливаются на максимум.
\end{itemize}

\booksubsection{Как включить простой режим на разных брендах (просто повторяйте действия):}

\begin{center}
\begin{tabularx}{\linewidth}{|X|X|X|}
\hline
\textbf{Бренд} & \textbf{Название режима} & \textbf{Как включить} \tabularnewline
\hline
iPhone & Универсальный доступ (отдельного простого режима нет, но можно увеличивать элементы по отдельности) & Настройки > Экран и яркость > Размер текста, вывести на максимум. Затем в «Универсальный доступ > Экран и размер текста» включить «Жирный текст» и «Формы кнопок» \tabularnewline
Huawei/Honor & Простой режим & Настройки > Система и обновление > Простой режим > Включить \tabularnewline
Xiaomi/Redmi & Минимальный режим & Настройки > Дополнительные настройки > Минимальный режим > Включить \tabularnewline
OPPO/OnePlus & Простой режим & Настройки > Специальные функции > Простой режим \tabularnewline
vivo/iQOO & Простой рабочий стол & Настройки > Рабочий стол, экран блокировки и обои > Простой рабочий стол \tabularnewline
Samsung & Простой режим & Настройки > Экран > Простой режим \tabularnewline
\hline
\end{tabularx}
\end{center}

Если у вас другой бренд, просто введите в поисковой строке вверху «Настроек» слова «простой» или «пожилой» — обычно функция находится. Если не нашли, настраивайте вручную по пунктам ниже, эффект будет тем же.

\booksection{Два. Если не включать простой режим, можно настроить вручную — «крупный шрифт + крупные значки» за пять шагов}

Некоторые пожилые не хотят использовать простой режим — им кажется, что интерфейс слишком «старомодный». Тогда вручную настройте следующие пять пунктов, эффект будет таким же:

\begin{enumerate}
  \item \textbf{Шрифт на максимум}: Настройки > Экран > Размер шрифта, перетащите ползунок вправо до конца. Заодно увеличьте «Размер экрана» (иногда называется «Масштаб экрана»), тогда все элементы интерфейса станут крупнее.
  \item \textbf{Значки крупнее}: На главном экране сведите два пальца (как при уменьшении фотографии) > Настройки рабочего стола > Макет (или размер значков), выберите макет с наименьшим числом ячеек (например 3×5) — значки автоматически станут больше.
  \item \textbf{Включите «Текст с высоким контрастом»}: края текста станут резче, этот переключатель находится в «Универсальном доступе».
  \item \textbf{Отключите лишние уведомления}: Настройки > Уведомления, выключите все, кроме Телефона, Сообщений и WeChat. Тогда телефон не будет каждый раз издавать звук и не будет раздражать.
  \item \textbf{Добавьте часто используемые контакты на рабочий стол}: Нажмите и удерживайте пустое место на рабочем столе > Добавить виджет > Контакты, разместите быстрый набор супруга/супруги, детей, врача, экстренных контактов. Одно нажатие — и звонок идёт, не нужно искать в телефонной книге.
\end{enumerate}

\begin{tcbpractice}{Практика: помогите пожилым «обставить» телефон}
\begin{enumerate}
  \item Включите простой режим (действуйте по таблице брендов выше).
  \item Включите отображение заряда в процентах (Настройки > Батарея > Процент заряда), чтобы они всегда знали, сколько осталось.
  \item Откройте панель уведомлений и уберите или скройте быстрые переключатели Bluetooth, NFC, авиарежима — пожилым они не нужны.
  \item Разместите на рабочем столе четыре кнопки быстрого набора: супруг/супруга, дети, семейный врач (если есть), управляющая компания / ЖКХ.
\end{enumerate}

\end{tcbpractice}

\booksection{Три. Защита от мошенничества — эту главу обязательно прочитайте вместе с пожилыми родственниками}

Пожилые люди — главная мишень телефонного мошенничества. Дело не в их наивности, а в том, что сценарии мошенников тщательно разработаны и бьют точно по психологическим уязвимостям старшего поколения. Ниже — самые распространённые виды мошенничества и способы защиты. Рекомендую зачитать это пожилым родственникам или распечатать и приклеить рядом с телефоном.

\booksubsection{Шесть самых частых мошеннических схем (если услышите — сразу кладите трубку):}

\begin{itemize}
  \item \textbf{(1) «Ваш внук попал в беду, срочно переведите деньги»}— мошенник говорит: «Ваш внук/внучка натворил дел, полиция его задержала, нужен залог». Запомните:\textbf{любой звонок с требованием перевести деньги за «освобождение» — это мошенничество}. Правоохранительные органы не требуют денег по телефону.
  \item \textbf{(2) «Ваша медицинская карта заблокирована»}— в SMS пишут, что с медицинской картой проблемы, и просят перейти по ссылке для «активации». Запомните:\textbf{медицинские и социальные фонды не рассылают SMS со ссылками}. В такой ситуации звоните на 12333 (официальная горячая линия социального страхования) для проверки.
  \item \textbf{(3) «Высокая доходность, гарантированная прибыль»}— в группах WeChat и SMS рекомендуют различные «пенсионные инвестиции» и «государственные проекты поддержки» с обещанием высокой прибыли. Запомните:\textbf{если обещают годовую доходность выше 5\% с «гарантией капитала» — на 99,9\% это мошенничество}.
  \item \textbf{(4) «Я ваш начальник/старый друг, сменил номер»}— выдают себя за знакомого и просят в долг. Запомните:\textbf{перзвонить для проверки — лучшее решение}. Если говорят «не могу говорить по телефону» — скорее всего, мошенник.
  \item \textbf{(5) «Бесплатные яйца / лекция о здоровье»}— бесплатное офлайн-мероприятие заманивает пенсионеров на «лекцию о здоровье», а затем продают БАДы по заоблачным ценам. Запомните:\textbf{бесплатный сыр только в мышеловке, за каждым бесплатным мероприятием стоит бизнес}.
  \item \textbf{(6) «Ваш телефон заражён вирусом, нужен удалённый ремонт»}— мошенник просит установить приложение удалённого управления (например, Sunlogin, TeamViewer), затем удалённо управляет вашим телефоном и переводит деньги. Запомните:\textbf{никогда не позволяйте никому удалённо управлять вашим телефоном!}Производители телефонов, операторы связи и банки не звонят с просьбой установить программу удалённого управления.
\end{itemize}

\booksubsection{Три линии обороны: настроив их, вы заблокируете 90\% мошенничества:}

\begin{enumerate}
  \item \textbf{Первая: встроенная блокировка телефона}— на Android включите «Блокировку спама» (в приложении Телефон или в Менеджере телефона), на iPhone включите «Заглушение неизвестных абонентов» (Настройки > Телефон > Заглушение неизвестных абонентов).
  \item \textbf{Вторая: блокировка оператором}— отправьте SMS «KTFSR» на 10086, чтобы бесплатно включить блокировку частых спам-звонков у China Mobile; у China Telecom и China Unicom есть аналогичные услуги. Просто позвоните в службу поддержки оператора и попросите включить.
  \item \textbf{Третья: приложение Национального центра противодействия мошенничеству}— скачайте приложение «Национальный центр противодействия мошенничеству» в магазине приложений, зарегистрируйтесь и включите предупреждение о входящих вызовах. Это приложение связано с базой данных МВД — как только звонит мошеннический номер, он определяется.
\end{enumerate}

\begin{tcbwarn}{«Правила защиты от мошенничества» для пожилых (распечатайте и наклейте рядом с телефоном)}
\begin{itemize}
  \item \textbf{Незнакомец просит перевести деньги кладите трубку!}
  \item \textbf{В SMS ссылка, которую нужно открыть удалите!}
  \item \textbf{Обещают высокую доходность с гарантией в чёрный список!}
  \item \textbf{Просят установить программу удалённого управления сразу сбросьте!}
  \item \textbf{Настоящая полиция не ведёт дел по телефону и не требует денег!}
  \item \textbf{Если сомневаетесь — сначала позвоните детям и спросите.}
\end{itemize}

\end{tcbwarn}

\booksection{Четыре. Удалённая помощь — даже из другого города вы можете починить телефон родителей}

Иногда проблему пожилых невозможно решить по телефону — вы полчаса объясняете, а они не могут найти нужную кнопку. Тут на помощь приходит функция «удалённой помощи» — вы можете напрямую видеть экран их телефона на своём и помогать с управлением. У каждого бренда свой способ:

\begin{itemize}
  \item \textbf{Huawei/Honor}: встроенное приложение «Changlian» (Canghai), в нём есть функция «Общий экран». Вы и родители устанавливаете приложение, начинаете видеозвонок и нажимаете «Общий экран» — вы можете обводить кружками элементы на их экране, показывая маме: «нажми сюда».
  \item \textbf{Xiaomi/Redmi}: в приложении «Видеозвонки Xiaomi» есть «Удалённая помощь» с общим экраном и синхронной голосовой связью.
  \item \textbf{OPPO/vivo}: встроенные «Семейное пространство» или «Удалённая опека», найдите в настройках по слову «удалённая».
  \item \textbf{Универсальный способ (подходит для всех брендов)}: скачайте приложение «Sunlogin Remote Control» (бесплатная версия достаточна), установите на оба телефона. Зарегистрируйте один аккаунт или добавьте друг друга, пусть родители продиктуют отображаемый на экране код подтверждения — и вы сможете удалённо видеть и управлять их экраном.
  \item \textbf{Для пользователей iPhone}: во время звонка FaceTime поддерживается «Общий экран». Нажмите кнопку демонстрации экрана, и собеседник увидит ваш экран. Но из-за ограничений безопасности iOS удалённое управление невозможно — только «инструктаж голосом».
\end{itemize}

\booksection{Пять. «Адаптация WeChat для пожилых» — достаточно сказать им эти несколько вещей}

WeChat — самое популярное приложение у пожилых людей, но интерфейс слишком перегружен. Вот несколько настроек с немедленным эффектом:

\begin{itemize}
  \item \textbf{Увеличьте шрифт в самом WeChat}: WeChat > Я > Настройки > Основные > Размер шрифта, выведите на максимум. Обратите внимание: это влияет только на шрифт внутри WeChat, не на системный шрифт.
  \item \textbf{Включите «Режим заботы»}: WeChat > Я > Настройки > Режим заботы > Включить. После включения шрифт и кнопки становятся крупнее, контрастность цветов усиливается, появляется возможность одним нажатием «слушать текстовые сообщения» (нажмите на текстовое сообщение — оно будет прочитано вслух).
  \item \textbf{Помогите подписать важные контакты}: измените имена детей и родственников на легко узнаваемые — например, «Старший сын - Сяомин», «Младшая дочь - Сяохун», чтобы при открытии WeChat они сразу понимали, кто есть кто.
  \item \textbf{Научите их использовать голосовой ввод вместо набора текста}: на клавиатуре рядом с полем чата удерживайте пробел (Android) или нажмите на значок микрофона (iPhone) — говорите, и речь автоматически превращается в текст. Точность распознавания очень высокая, для пожилых людей, не знающих пиньинь, это огромное облегчение.
  \item \textbf{Отключите неиспользуемые функции во вкладке «Обнаружить»}: WeChat > Я > Настройки > Основные > Управление вкладкой «Обнаружить», выключите все ненужные: «Моменты», «Каналы», «Трансляции», «Мини-шоп», «Игры» — меньше отвлечений.
\end{itemize}

\booksection{Шесть. Что делать, если телефон потерялся? Помогите пожилым составить «аварийный план»}

Предусмотрительность — залог успеха. Пока телефон на месте, сделайте следующие приготовления:

\begin{enumerate}
  \item \textbf{Убедитесь, что «Найти устройство» включено}: на iPhone — Настройки > Ваше имя > Найти > Найти iPhone; на Android — в настройках найдите «Найти устройство». Обязательно проверьте, что переключатель включён.
  \item \textbf{Запишите важные пароли и аккаунты на листе бумаги}: пароли от Apple ID / аккаунта Huawei / аккаунта Xiaomi. Один экземпляр — в запертом ящике дома, другой — у доверенного ребёнка. После потери телефона понадобится войти в эти аккаунты для удалённой блокировки и отслеживания; если забудете пароль — будет очень проблемно.
  \item \textbf{Научите пожилых одной фразе}:\textbf{«Телефон потерялся — первым делом позвони семье».}Родные смогут немедленно зайти на iCloud или сайт «Найти устройство» для Android и удалённо включить «режим пропажи» — тогда нашедший не сможет ни открыть, ни использовать телефон.
\end{enumerate}

\booksection{Сем. Совместимость беспроводных наушников и слуховых аппаратов — недооценённая полезная функция}

Многие пожилые имеют проблемы со слухом, но не знают, что телефон может работать со слуховыми аппаратами. iPhone и некоторые Android-телефоны поддерживают слуховые аппараты MFi (Made for iPhone) и могут передавать звук звонка прямо в слуховой аппарат — гораздо чётче, чем прикладывая телефон к уху.

\begin{itemize}
  \item \textbf{Как подключить}: Настройки iPhone > Универсальный доступ > Слуховые устройства, включите Bluetooth и следуйте инструкциям по сопряжению слухового аппарата.
  \item \textbf{Если слуховой аппарат не используется}: в «Универсальный доступ > Аудио/визуальные» включите «Шумоподавление телефона» и «Моно-аудио» — это улучшит чёткость разговора.
  \item \textbf{Android}: в «Универсальный доступ > Усиление слуха» есть аналогичные настройки, расположение немного отличается у разных брендов — просто найдите слово «слух».
\end{itemize}

\booksection{Восемь. Функции здоровья — телефон не только средство связи}

\begin{itemize}
  \item \textbf{SOS-экстренный вызов}: пять раз подряд нажмите кнопку питания (на большинстве телефонов) — телефон автоматически позвонит в экстренные службы и отправит геолокацию заранее заданным контактам. Настройте это для пожилых заранее — iPhone: «Настройки > Экстренный вызов SOS», Android: «Настройки > Безопасность > Экстренный вызов». Перед тем как учить пожилых пользоваться, потренируйтесь вместе, чтобы в ответственный момент не нажимали с ошибкой от волнения.
  \item \textbf{Обнаружение падения}: эта функция есть в Apple Watch — если часы определяют, что вы упали и не двигаетесь какое-то время, они автоматически вызывают скорую помощь и уведомляют близких. Если бюджет позволяет, подарите пожилому родственнику такие часы — в критический момент они могут спасти жизнь.
  \item \textbf{Напоминание о приёме лекарств}: в стандартном приложении «Здоровье» на iPhone есть встроенная функция напоминаний о лекарствах — задайте название, время и дозировку, и в нужный момент появится уведомление. На Android можно использовать «таймер-будильник» или скачать специальные приложения-напоминалки.
\end{itemize}

\begin{tcbsummary}{Итоги главы}
\begin{itemize}
  \item Сначала включите простой режим или вручную увеличьте шрифт/иконки — после настройки среды большая часть проблем исчезнет
  \item Три железных правила обучения пожилых: не торопите, используйте аналогии, записывайте
  \item Шесть основных видов мошенничества нужно постоянно напоминать — три линии обороны отсеют 90\% мошенничества
  \item Удалённая помощь позволяет чинить телефон родителей, даже находясь в другом городе
  \item Включение «режима заботы» и «голосового ввода» в WeChat — две самые полезные функции для пожилых
  \item Экстренный вызов SOS и обнаружение падения — настройки, спасающие жизнь в критический момент, стоят пяти минут настройки
\end{itemize}

\end{tcbsummary}
